% Drawing 77/05-3 -- folder 77, batch 05, archive page 87 (PDF page 88).
% Transcription: deux rangées de tirets, et un hexagone.
% Kind: hexagon and two rows of dashes. Method 2 (schematic redraw from the facsimile). Status: ok.
% Notes: An open-stroked hexagon, and two rows of dashes to the right of and below it.
% Coordinates: 1 unit = 1 mm = 1/100 of the crop width (bbox in page fractions: [0.52, 0.34, 0.76, 0.43]).
% Redrawn by Claude (claude-opus-5-5), issue #73, 2026-11-07; not checked by a human; not compiled (no TeX in the session).
% Rights: reproduces part of the work -- behind RENDER_TRANSCRIPTS, never published by itself.
\documentclass[tikz,border=3pt]{standalone}
\usepackage{amsmath,amssymb}
\usetikzlibrary{patterns,arrows.meta}
\tikzset{pencil/.style={black!45}, struck/.style={black!35}, mark/.style={circle,draw,inner sep=1.3pt}, pt/.style={circle,fill,inner sep=0.9pt}, opt/.style={circle,draw,fill=white,inner sep=1.1pt}}
\begin{document}
\begin{tikzpicture}[x=1mm,y=1mm,line cap=round,line join=round,font=\scriptsize]
\draw (21,-22) -- (32,-11);
\draw (33,-12.5) -- (49,-13);
\draw (51,-14) -- (59,-24.5);
\draw (60,-27) -- (48,-40);
\draw (47,-40) -- (31,-39.5) -- (21,-24.5);
\draw (77,-29.5) -- (91,-29.5);
\draw (98,-29.8) -- (100,-29.8);
\draw (0,-52) -- (13,-52);
\draw (24,-52) -- (38,-52);
\draw (49,-52) -- (63,-52);
\draw (75,-52) -- (88,-52);
\draw (98,-52) -- (100,-52);
\end{tikzpicture}
\end{document}
